\chapter{Haemophilia}
\index{Haemophilia}

\section*{Definition and Overview}
Haemophilia is an inherited bleeding disorder in which the blood does not clot properly because of a deficiency of a specific clotting factor. Haemophilia A is caused by a lack of clotting factor VIII, and haemophilia B by a lack of factor IX. The condition is caused by a genetic change on the X chromosome, so it affects almost exclusively males, while females are usually carriers. People with haemophilia bleed for longer after injury, and can bleed spontaneously into joints and muscles, causing pain, swelling, and, over time, joint damage. Modern treatment replaces the missing clotting factor, allowing people with haemophilia to live active lives, and most manage with factor infusions at home.

\section*{Epidemiology}
Haemophilia is rare, affecting around 1 in 5,000--10,000 males for haemophilia A and about 1 in 30,000--50,000 for haemophilia B. In many countries, around 3,000 people live with haemophilia. It occurs in all ethnic groups. Severity depends on the level of clotting factor: around half of cases are severe (less than 1\% factor), and people with severe haemophilia bleed spontaneously. About a third of cases arise from new mutations, so they occur with no family history. Diagnosis is often made in infancy or early childhood, sometimes after circumcision or with bruising when learning to walk.

\section*{Aetiology and Pathophysiology}
Haemophilia is caused by genetic changes that reduce or eliminate a clotting factor.

\begin{itemize}
    \item \textbf{Genetic basis:} mutations in the F8 gene (haemophilia A) or F9 gene (haemophilia B), located on the X chromosome
    \item \textbf{X-linked inheritance:} males with the affected gene have the disease; females are usually carriers but can be affected in rare circumstances
    \item \textbf{Clotting cascade:} factor VIII and factor IX are essential for forming a stable blood clot; without them, clotting is delayed and bleeding continues
    \item \textbf{Severity:} factor levels of less than 1\% cause severe disease, 1--5\% moderate, and 5--40\% mild
    \item \textbf{Spontaneous bleeding:} in severe disease, bleeding occurs into joints and muscles without injury
    \item \textbf{Inhibitors:} a small proportion of people develop antibodies (inhibitors) that block the replacement factor
\end{itemize}

\section*{Clinical Features}
Symptoms depend on the severity of the deficiency.

\begin{itemize}
    \item Prolonged bleeding after cuts, dental work, or surgery
    \item Easy bruising, sometimes with large bruises
    \item Bleeding into joints (haemarthrosis), causing pain, swelling, and stiffness, most often in knees, elbows, and ankles
    \item Bleeding into muscles, causing swelling and pain
    \item Spontaneous bleeding in severe haemophilia, without injury
    \item Prolonged bleeding after circumcision, often the first sign in newborns
    \item Nosebleeds and mouth bleeding that are hard to stop
    \item Blood in the urine or stool
    \item Bleeding into the brain (intracranial haemorrhage), which is rare but serious, especially after head injury
\end{itemize}

\section*{Differential Diagnosis}
Other bleeding disorders should be considered.

\begin{itemize}
    \item von Willebrand disease, the most common inherited bleeding disorder
    \item Other clotting factor deficiencies, including factor XI deficiency
    \item Platelet disorders and platelet function defects
    \item Vitamin K deficiency in newborns
    \item Acquired bleeding disorders, including those from liver disease or medicines
    \item Acquired inhibitors of clotting factors
    \item Non-accidental injury, which must be considered when bleeding is unexplained
\end{itemize}

\section*{When to Seek Medical Care}
Families with a history of haemophilia should have genetic counselling and testing. Anyone with unusual or prolonged bleeding should be assessed by a doctor, and people with known haemophilia should have a treatment plan and know when to seek care.

\textbf{Contact your local emergency services for:}
\begin{itemize}
    \item Head injury in a person with haemophilia, which requires urgent assessment
    \item Sudden severe headache, vomiting, or neurological symptoms, which may indicate intracranial bleeding
    \item Bleeding that will not stop despite factor treatment
    \item Severe joint or muscle bleeding causing significant pain and swelling
    \item Blood in the urine or stool, or any heavy bleeding
\end{itemize}

\section*{Diagnosis and Investigations}
Haemophilia is diagnosed by clotting tests and factor assays.

\begin{itemize}
    \item \textbf{Clotting tests:} prolonged activated partial thromboplastin time (APTT) with normal prothrombin time
    \item \textbf{Factor assays:} measurement of factor VIII and factor IX levels, which also determines severity
    \item \textbf{Genetic testing:} identifies the specific mutation and supports carrier testing and prenatal diagnosis
    \item \textbf{Carrier testing:} offered to female relatives
    \item \textbf{Prenatal testing:} available for families with known mutations
    \item \textbf{Newborn screening:} not routine, but testing is arranged when there is a family history or unusual bleeding
    \item \textbf{Inhibitor testing:} if factor treatment stops working
\end{itemize}

\section*{Management}
Modern management allows people with haemophilia to live full and active lives.

\begin{itemize}
    \item \textbf{Factor replacement:} intravenous infusions of factor VIII or IX, given to treat bleeding and, in prophylaxis, to prevent it
    \item \textbf{Prophylaxis:} regular factor infusions, usually from childhood, which prevent bleeding and joint damage
    \item \textbf{Home treatment:} most people are trained to infuse factor at home
    \item \textbf{Treatment of bleeding:} prompt factor infusion at the first sign of bleeding
    \item \textbf{Inhibitor management:} immune tolerance therapy and alternative agents for people with inhibitors
    \item \textbf{Non-factor therapies:} newer medicines that promote clotting by other mechanisms
    \item \textbf{Physiotherapy:} rehabilitation of joints and muscles after bleeding
    \item \textbf{Comprehensive care:} management at a haemophilia treatment centre with a multidisciplinary team
    \item \textbf{Surgical and dental care:} factor cover before procedures
\end{itemize}

\section*{Complications}
\begin{itemize}
    \item Recurrent joint bleeding leading to haemophilic arthropathy, chronic pain, and disability
    \item Muscle bleeding causing nerve compression
    \item Intracranial bleeding, which can be fatal or cause brain injury
    \item Development of inhibitors, making treatment more difficult
    \item Transfusion-transmitted infections, now very rare with modern products
    \item Anaemia from chronic bleeding
    \item Complications of treatment, including allergic reactions
    \item Effects on schooling, work, and quality of life
\end{itemize}

\section*{Prognosis and Outcomes}
The outlook for people with haemophilia has transformed over recent decades. With prophylactic factor replacement, children with severe haemophilia grow up without the crippling joint disease of the past, and people with mild haemophilia may rarely need treatment. Life expectancy for people with haemophilia in developed countries now approaches that of the general population. New therapies, including longer-acting factors and non-factor medicines, continue to improve outcomes. With comprehensive care, people with haemophilia can participate in most activities, with sensible precautions.

\section*{Prevention and Self-Care}
\begin{itemize}
    \item Attend a comprehensive haemophilia treatment centre
    \item Follow the prescribed prophylaxis schedule
    \item Treat bleeding early, at the first sign
    \item Maintain a healthy weight to protect the joints
    \item Choose safe physical activities, and wear protective gear for contact sports
    \item Avoid medicines that worsen bleeding, including aspirin and other NSAIDs
    \item Practise good dental hygiene to reduce dental bleeding
    \item Keep immunisations up to date, using appropriate injection techniques
    \item Carry a medical alert card or wear a bracelet
    \item Seek genetic counselling and testing for family planning
\end{itemize}

\section*{Sources and Further Reading}
The following reputable sources provide further information for healthcare students and patients seeking reliable, current advice about haemophilia.

\begin{itemize}
    \item Haemophilia Foundation Australia. \textit{Information and support.} https://www.haemophilia.org.au/
    \item Australian Haemophilia Centre Directors' Organisation. \textit{Treatment centre directory and resources.} https://www.ahcdo.org.au/
    \item Healthdirect Australia. \textit{Haemophilia.} https://www.healthdirect.gov.au/
    \item World Federation of Hemophilia. \textit{Information about haemophilia.} https://wfh.org/
    \item NHS. \textit{Haemophilia.} https://www.nhs.uk/
    \item Centers for Disease Control and Prevention. \textit{Hemophilia.} https://www.cdc.gov/
\end{itemize}
